\chapter{Expanding to $n$-exact categories}\label{chapter:nexact}
After generalizing from abelian categories to exact categories, we can go one step further by looking at longer exact sequences. First, we will redefine the notions needed for exact categories, but in a more general context. Then we will proceed to define $n$-exact categories, and give some concrete example of such categories.

When not stated otherwise, $\cat{A}$ denotes an additive category. The contents of this chapter are based on \cite{Jasso2016Aug}.

\section{Working in higher degrees}
Before defining $n$-exact categories, there are many prerequisites that need to be addressed. Firstly, we will need a way to properly describe what the equivalent of a kernel-cokernel pair would be in this new context. As a result, it makes sense to define $n$-kernels and $n$-cokernels, which require the notion of weak kernel (respectively weak cokernel).

\begin{defn}
    Let $A \xra{f} B$ be a morphism in $\cat{A}$. A \emph{weak kernel} of $f$ is an object $K \in \cat{A}$ with a morphism $K \xra{k} A$ such that the sequence of abelian groups

    % https://q.uiver.app/#q=WzAsMyxbMCwwLCJcXGNhdHtBfShDJyxLKSJdLFsyLDAsIlxcY2F0e0F9KEMnLEEpIl0sWzQsMCwiXFxjYXR7QX0oQycsQikiXSxbMCwxLCItIFxcY2RvdCBrIl0sWzEsMiwiLSBcXGNkb3QgZiJdXQ==
    \[\begin{tikzcd}[ampersand replacement=\&,sep=scriptsize]
        {\cat{A}(C',K)} \&\& {\cat{A}(C',A)} \&\& {\cat{A}(C',B)}
        \arrow["{k \circ -}", from=1-1, to=1-3]
        \arrow["{f \circ -}", from=1-3, to=1-5]
    \end{tikzcd}\]

    is exact for all $C'\in \cat{A}$.

    Dually, a \emph{weak cokernel} of $f$ is an object $C \in \cat{A}$ with a morphism $B \xra{g} C$ such that the sequence of abelian groups 
    
    \[\begin{tikzcd}[ampersand replacement=\&,sep=scriptsize]
        {\cat{A}(C,C')} \&\& {\cat{A}(B,C')} \&\& {\cat{A}(A,C')}
        \arrow["{- \circ k}", from=1-1, to=1-3]
        \arrow["{- \circ f}", from=1-3, to=1-5]
    \end{tikzcd}\]

    is exact for all $C' \in \cat{A}$.
\end{defn}

\begin{rmks} We do not ask for the sequence to be short exact, but only for it to be exact in the middle term. Furthermore, the following observations are quite useful in practice:
    \begin{enumerate}[label = (\alph*)]
        \item Notice the direction of the sequences is different since for weak cokernels, we apply the contravariant functor $\cat{A}(-, C')$ to our original sequence, whereas we apply the covariant functor $\cat{A}(C', -)$ in the case of weak kernels.
        \item For $g: K \rightarrow A$ a weak kernel of $f: A \rightarrow B$, the sequence being exact lets us know that
        \begin{enumerate}
            \item The composition $f \cdot g$ is $0$;
            \item For any morphism $h: C \rightarrow A$ such that \break \mbox{$f \cdot h = 0$}, there exists a morphism $\hat{h}: C \xra{} K$ such that $g \circ \hat{h} = h$.
        \end{enumerate}

        That last point can be summarized by the following diagram:

        % https://q.uiver.app/#q=WzAsNCxbMCwwLCJLIl0sWzIsMCwiQSJdLFs0LDAsIkIiXSxbMCwyLCJDIl0sWzMsMSwiaCJdLFsxLDIsImYiXSxbMCwxLCJnIl0sWzMsMiwiMCIsMix7Im9mZnNldCI6MywiY3VydmUiOi0zfV0sWzMsMCwiXFxleGlzdHMgXFxoYXR7aH0iLDAseyJzdHlsZSI6eyJib2R5Ijp7Im5hbWUiOiJkYXNoZWQifX19XV0=
        \[\begin{tikzcd}[ampersand replacement=\&,sep=scriptsize]
            K \&\& A \&\& B \\
            \\
            C
            \arrow["g", from=1-1, to=1-3]
            \arrow["f", from=1-3, to=1-5]
            \arrow["{\exists \hat{h}}", dashed, from=3-1, to=1-1]
            \arrow["h", from=3-1, to=1-3]
            \arrow["0"', shift right=2, curve={height=-15pt}, from=3-1, to=1-5]
        \end{tikzcd}\]

        These two points follow from exactness in the middle term. Indeed, since $\Img{(k \circ -)} = \Ker{(f \circ -)}$, any map $h: C \rightarrow A$ such that $f \circ h = 0$, is in the image of the post-composition with $k$, i.e. can be factored (not necessarily uniquely!) as $k \circ \hat{h}$. In particular, if we consider $C' = K$, we conclude that 
        \[k = k \circ 1_K \in \Img{(k \circ -)} = \Ker{(f \circ -)},\]
        i.e. $f \circ k = 0$.
        The result for weak cokernels is dual.
    \end{enumerate}

    In summary, weak cokernels and weak kernels have a very similar property to that of cokernels and kernels, except that uniqueness in our factoring of morphisms is not guaranteed. In other words, the only difference is that a weak kernel is not necessarily monic, and a weak cokernel is not necessarily epic.
\end{rmks}

This notion is quite fundamental in the elaboration of $n$-kernels and $n$-cokernels:

\begin{defn}
    Let $X^0 \xra{d^0} X^1$ be a morphism in $\cat{A}$. An \emph{$n$-kernel} of $d^0$ is a sequence 

    % https://q.uiver.app/#q=WzAsNCxbMCwwLCJLXjAiXSxbMiwwLCJLXjEiXSxbNiwwLCJLXm4iXSxbNCwwLCJcXGxkb3RzIl0sWzAsMV0sWzEsM10sWzMsMl1d
    \[\begin{tikzcd}[ampersand replacement=\&,sep=scriptsize]
        {X^{-n}} \&\& {X^{-n+1}} \&\& \cdots \&\& {X^0}
        \arrow["d^{-n}",from=1-1, to=1-3]
        \arrow["d^{-n+1}", from=1-3, to=1-5]
        \arrow["d^{0}", from=1-5, to=1-7]
    \end{tikzcd}\]

    such that the sequence induced by applying the functor $\cat{A}(C', -)$
    % https://q.uiver.app/#q=WzAsNSxbMiwwLCJcXGNhdHtBfShDJywgWF57LW59KSJdLFs0LDAsIlxcY2F0e0F9KEMnLCBYXnstbisxfSkiXSxbOCwwLCJcXGNhdHtBfShDJywgWF57MH0pIl0sWzYsMCwiXFxsZG90cyJdLFswLDAsIjAiXSxbMCwxLCJkXnstbn0gXFxjaXJjIC0iXSxbMSwzLCJkXnstbisxfSBcXGNpcmMgLSJdLFszLDIsImReey0xfSBcXGNpcmMgLSJdLFs0LDBdXQ==
    \[\begin{tikzcd}[ampersand replacement=\&,sep=small]
        0 \&\& {\cat{A}(C', X^{-n})} \&\& {\cat{A}(C', X^{-n+1})} \&\& \ldots \&\& {\cat{A}(C', X^{0})}
        \arrow[from=1-1, to=1-3]
        \arrow["{d^{-n} \circ -}", from=1-3, to=1-5]
        \arrow["{d^{-n+1} \circ -}", from=1-5, to=1-7]
        \arrow["{d^{-1} \circ -}", from=1-7, to=1-9]
    \end{tikzcd}\]
    is exact.

    Equivalently, an $n$-kernel of $d^0$ is such a sequence $(d^{-n}, d^{-n+1}, \ldots, d^{-1})$, where $d^{-k}$ is a weak cokernel of $d^{-k+1}$ for all $1 \leq k \leq n-1$, and $d^{-n}$ is a kernel of $d^{-n+1}$.

    Dually, an \emph{$n$-cokernel} of $d^0$ is a sequence

    % https://q.uiver.app/#q=WzAsNCxbMCwwLCJYXjEiXSxbMiwwLCJYXjIiXSxbNCwwLCJcXGxkb3RzIl0sWzYsMCwiWF57bisxfSJdLFswLDEsImReMSJdLFsxLDIsImReMiJdLFsyLDMsImRebiJdXQ==
    \[\begin{tikzcd}[ampersand replacement=\&,sep=small]
        {X^1} \&\& {X^2} \&\& \cdots \&\& {X^{n+1}}
        \arrow["{d^1}", from=1-1, to=1-3]
        \arrow["{d^2}", from=1-3, to=1-5]
        \arrow["{d^n}", from=1-5, to=1-7]
    \end{tikzcd}\]

    such that $d^{k}$ is a weak cokernel of $d^{k-1}$ for $1 \leq k \leq n-1$ and $d^n$ is a cokernel of $d^{n-1}$.
\end{defn}

\begin{rmk}
    In the case $n=1$, this indeed coincides with the usual definition of kernels and cokernels.
\end{rmk}

Now that we have a suitable notion of higher kernels and cokernels, we can consider the concept of $n$-exact sequences in the most logical way.

\begin{defn}
    An \emph{$n$-exact sequence} in $\cat{A}$ is a chain complex 
    \[X^0 \xra{d^0} X^1 \xra{d^1} \cdots \xra{d^{n-1}} X^n \xra{d^n} X^{n+1}\]
    such that $(d^0, \ldots, d^{n-1})$ is an $n$-kernel of $d^n$ and $(d^1, \ldots, d^n)$ is an ${n\text{-cokernel}}$ of $d^0$.
\end{defn}

\begin{rmk}
    Once again, when $n=1$, we get back to our usual notion of short exact sequence.
\end{rmk}

There are two final notions that need to be covered before proceeding with the definition of $n$-exact categories: that of $n$-pullback, $n$-pushout, and weak isomorphisms. Before covering these concepts, let us first remember what the cone of a morphism of complexes is:

\begin{defn}
    \begin{sloppypar}
    Let $X$ and $Y$ be two complexes in $\Ch{\cat{A}}$ and ${f: X \rightarrow Y}$ be a morphism of complexes, or diagrammatically:
    \end{sloppypar}

    % https://q.uiver.app/#q=WzAsMTIsWzAsMCwiWCJdLFswLDIsIlkiXSxbMiwwLCJcXGxkb3RzIl0sWzQsMCwiWF57bi0xfSJdLFs2LDAsIlhebiJdLFsyLDIsIlxcbGRvdHMiXSxbNCwyLCJZXntuLTF9Il0sWzYsMiwiWV5uIl0sWzgsMCwiWF57bisxfSJdLFs4LDIsIllee24rMX0iXSxbMTAsMCwiXFxsZG90cyJdLFsxMCwyLCJcXGxkb3RzIl0sWzAsMSwiZiJdLFsyLDNdLFszLDQsImRfWF57bi0xfSJdLFs1LDZdLFs2LDcsImRfWV57bi0xfSJdLFs0LDgsImRfWF5uIl0sWzcsOSwiZF9ZXntufSJdLFs4LDEwXSxbOSwxMV0sWzgsOSwiZl57bisxfSJdLFs0LDcsImZebiJdLFszLDYsImZee24tMX0iXV0=
    \[\begin{tikzcd}[ampersand replacement=\&,sep=small]
        X: \&\& \cdots \&\& {X^{n-1}} \&\& {X^n} \&\& {X^{n+1}} \&\& \cdots \\
        \\
        Y: \&\& \cdots \&\& {Y^{n-1}} \&\& {Y^n} \&\& {Y^{n+1}} \&\& \cdots
        \arrow["f", from=1-1, to=3-1]
        \arrow[from=1-3, to=1-5]
        \arrow["{d_X^{n-1}}", from=1-5, to=1-7]
        \arrow["{f^{n-1}}", from=1-5, to=3-5]
        \arrow["{d_X^n}", from=1-7, to=1-9]
        \arrow["{f^n}", from=1-7, to=3-7]
        \arrow[from=1-9, to=1-11]
        \arrow["{f^{n+1}}", from=1-9, to=3-9]
        \arrow[from=3-3, to=3-5]
        \arrow["{d_Y^{n-1}}", from=3-5, to=3-7]
        \arrow["{d_Y^{n}}", from=3-7, to=3-9]
        \arrow[from=3-9, to=3-11]
    \end{tikzcd}\]

    The \emph{cone} of $f$, denoted $\C{f}$, is the following complex

    % https://q.uiver.app/#q=WzAsNSxbMCwwLCJDKGYpOiJdLFsxLDAsIlxcY2RvdHMiXSxbMywwLCJYXm4gXFxiaWdvcGx1cyBZXntuLTF9Il0sWzUsMCwiWF57bisxfSBcXGJpZ29wbHVzIFlee259Il0sWzcsMCwiXFxjZG90cyJdLFsxLDJdLFsyLDMsImRfQ157bi0xfSJdLFszLDRdXQ==
    \[\begin{tikzcd}[ampersand replacement=\&,sep=small]
        {C(f):} \& \cdots \&\& {X^n \bigoplus Y^{n-1}} \&\& {X^{n+1} \bigoplus Y^{n}} \&\& \cdots
        \arrow[from=1-2, to=1-4]
        \arrow["{d_C^{n-1}}", from=1-4, to=1-6]
        \arrow[from=1-6, to=1-8]
    \end{tikzcd}\]

    where $d_C^{n-1} = \begin{pmatrix}-d_X^{n} & 0 \\ f^n & d_Y^{n-1}\end{pmatrix}$.
\end{defn}

Before looking at the higher equivalents of pullbacks and pushouts, it is quite useful to examine the link that they have with the cone of a map.

\begin{rmk}
    Let $B^1 \xra{d_B^1} B^2$ and $A^2 \xra{f^2} B^2$ be two morphisms in $\cat{A}$ and consider the pullback of $d_B^1$ along $f^2$:

    \[\begin{tikzcd}[ampersand replacement=\&, sep=small]
        {A^1} \&\& {A^2} \\
        \& {\text{PB}} \\
        B^1 \&\& B^2
        \arrow["{d_A^{1}}", shorten <=5pt, shorten >=5pt, from=1-1, to=1-3]
        \arrow["\lowercase{f^1}"', shorten <=3pt, shorten >=3pt, from=1-1, to=3-1]
        \arrow["\lowercase{f^2}", shorten <=3pt, shorten >=3pt, from=1-3, to=3-3]
        \arrow["d_B^1"', shorten <=5pt, shorten >=5pt, from=3-1, to=3-3]
    \end{tikzcd}\]

    Clearly $f = (f^1, f^2)$ is a morphism of complexes since the square commutes. Thus we can consider its cone:

    % https://q.uiver.app/#q=WzAsNixbMCwwLCJDKGYpOiJdLFsxLDAsIjAiXSxbMywwLCJBXjEiXSxbNSwwLCJBXjIgXFxiaWdvcGx1cyBCXjEiXSxbNywwLCJCXjIiXSxbOSwwLCIwIl0sWzEsMl0sWzIsMywiZF9DXjEiXSxbMyw0LCJkX0NeMiJdLFs0LDVdXQ==
    \[\begin{tikzcd}[ampersand replacement=\&,sep=small]
        {\C{f}:} \& 0 \&\& {A^1} \&\& {A^2 \bigoplus B^1} \&\& {B^2} \&\& 0
        \arrow[from=1-2, to=1-4]
        \arrow["{d_C^1}", from=1-4, to=1-6]
        \arrow["{d_C^2}", from=1-6, to=1-8]
        \arrow[from=1-8, to=1-10]
    \end{tikzcd}\]

    where $d_C^1 = \begin{pmatrix} - d_A^1 \\ f^1 \end{pmatrix}$, and $d_C^2 = \begin{pmatrix} f^2 & d_B^1 \end{pmatrix}$. We claim that $d_C^1$ is a kernel of $d_C^2$. Indeed, consider a morphism $\begin{pmatrix} g \\ h \end{pmatrix}: C \rightarrow A^2 \bigoplus B^1$ such that the composition $ d_C^2 \circ \begin{pmatrix} g \\ h \end{pmatrix}= 0$. The composition being $0$ implies that $ f^2 \circ g + d_B^1 \circ h = 0$, or $d_B^1 \circ h = - f^2 \circ g$. Consider the unique map $\hat{h} : C \rightarrow A^1$ given by $f^1 \circ \hat{h} = h$ and $d_A^1 \circ \hat{h} = -g$ (this map comes from the pullback property). We see that $d_C^1 \circ \hat{h} = \begin{pmatrix} - d_A^1 \circ \hat{h} \\ f^1 \circ \hat{h} \end{pmatrix} = \begin{pmatrix} g \\ h \end{pmatrix}$. Thus we have defined a unique map $\hat{h} : C \rightarrow A^1$ such that $d_C^1 \circ \hat{h} = \begin{pmatrix} g \\ h \end{pmatrix}$. That is, $d_C^1$ is a kernel of $d_C^2$.

    A dual result is true for pushouts, where $d_C^2$ is a cokernel of $d_C^1$.
\end{rmk}

Basing ourselves on this observation, we can generalize the notion of pullback and pushout to higher degrees:

\begin{defn}
    Let $X : X^0 \xra{d_X^0} X^1 \xra{d_X^1} \ldots \xra{d_X^{n-1}} X^n$ be a complex in $\Chlen{n}{\cat{A}}$ and $f^0 : X^0 \rightarrow Y^0$ a morphism in $\cat{A}$. An \emph{$n$-pushout of $X$ along $f^0$} is a morphism of complexes 

    % https://q.uiver.app/#q=WzAsMTAsWzAsMCwiWDoiXSxbMCwyLCJZOiJdLFsyLDAsIlheMCJdLFsyLDIsIlleMCJdLFs0LDAsIlheMSJdLFs0LDIsIlleMSJdLFs2LDAsIlxcY2RvdHMiXSxbNiwyLCJcXGNkb3RzIl0sWzgsMCwiWF5uIl0sWzgsMiwiWV5uIl0sWzAsMSwiZiJdLFsyLDQsImRfWF4wIl0sWzMsNSwiZF9ZXnswfSIsMCx7InN0eWxlIjp7ImJvZHkiOnsibmFtZSI6ImRhc2hlZCJ9fX1dLFsyLDMsImZeMCJdLFs0LDUsImZeMSIsMCx7InN0eWxlIjp7ImJvZHkiOnsibmFtZSI6ImRhc2hlZCJ9fX1dLFs0LDYsImRfWF4xIl0sWzUsNywiZF9ZXnsxfSIsMCx7InN0eWxlIjp7ImJvZHkiOnsibmFtZSI6ImRhc2hlZCJ9fX1dLFs2LDgsImRfWF57bi0xfSJdLFs3LDksImRfWV57bi0xfSIsMCx7InN0eWxlIjp7ImJvZHkiOnsibmFtZSI6ImRhc2hlZCJ9fX1dLFs4LDksImZebiIsMCx7InN0eWxlIjp7ImJvZHkiOnsibmFtZSI6ImRhc2hlZCJ9fX1dXQ==
    \[\begin{tikzcd}[ampersand replacement=\&,sep=small]
        {X:} \&\& {X^0} \&\& {X^1} \&\& \cdots \&\& {X^n} \\
        \\
        {Y:} \&\& {Y^0} \&\& {Y^1} \&\& \cdots \&\& {Y^n}
        \arrow["f", from=1-1, to=3-1]
        \arrow["{d_X^0}", from=1-3, to=1-5]
        \arrow["{f^0}", from=1-3, to=3-3]
        \arrow["{d_X^1}", from=1-5, to=1-7]
        \arrow["{f^1}", dashed, from=1-5, to=3-5]
        \arrow["{d_X^{n-1}}", from=1-7, to=1-9]
        \arrow["{f^n}", dashed, from=1-9, to=3-9]
        \arrow["{d_Y^{0}}", dashed, from=3-3, to=3-5]
        \arrow["{d_Y^{1}}", dashed, from=3-5, to=3-7]
        \arrow["{d_Y^{n-1}}", dashed, from=3-7, to=3-9]
    \end{tikzcd}\]

    such that in the cone of $f$
    % https://q.uiver.app/#q=WzAsNixbMSwwLCJYXjAiXSxbMiwwLCJYXjEgXFxiaWdvcGx1cyBZXjAiXSxbMywwLCJcXGxkb3RzIl0sWzQsMCwiWF57bn0gXFxiaWdvcGx1cyBZXntuLTF9Il0sWzUsMCwiWV5uIl0sWzAsMCwiXFxDe2Z9OiJdLFswLDEsImRfQ157LTF9Il0sWzEsMiwiZF9DXjAiXSxbMiwzLCJkX0Nee24tMX0iXSxbMyw0LCJkX0Nee259Il1d
    \[\begin{tikzcd}[ampersand replacement=\&,sep=scriptsize]
        {\C{f}:} \& {X^0} \& {X^1 \bigoplus Y^0} \& \cdots \& {X^{n} \bigoplus Y^{n-1}} \& {Y^n}
        \arrow["{d_C^{-1}}", from=1-2, to=1-3]
        \arrow["{d_C^0}", from=1-3, to=1-4]
        \arrow["{d_C^{n-1}}", from=1-4, to=1-5]
        \arrow["{d_C^{n}}", from=1-5, to=1-6]
    \end{tikzcd}\]

    $(d_C^{0}, \ldots, d_C^{n})$ is an $n$-cokernel of $d_C^{-1}$.

    The definition for the $n$-pullback of a chain complex is dual. That is, given a complex $Y$ in $\Chlen{n}{\cat{A}}$ and a morphism $f^n: X^n \rightarrow Y^n$, an \emph{$n$-pullback of $Y$ along $f^n$} is a morphism of complexes $f: X \rightarrow Y$ such that in its cone $C = \C{f}$, the sequence $(d_C^{-1}, \ldots, d_C^{n-1})$ is an $n$-kernel of $d_C^n$.
\end{defn}

The last notion required for defining $n$-exact categories is that of weak isomorphisms:

\begin{defn}
    Let $X$ and $Y$ be two $n$-exact sequences in $\cat{A}$. We say $f: X \rightarrow Y$ is a \emph{weak isomorphism} of complexes if there is $k \in \left\{0, 1, \ldots, n\right\}$ such that $f^k$ and $f^{k+1}$ are isomorphisms.
\end{defn}

When $n = 1$, we note that weak isomorphisms are isomorphisms of short exact sequences. This can be shown using the five lemma, which holds for any exact category (see \cite[Lemma~8.9]{Buhler2010Jan}) In the general case, however, weak isomorphisms are not necessarily isomorphisms, but a weaker result still holds:

\begin{prop}
    Let $X$ and $Y$ be two $n$-exact sequences in $\cat{A}$ and let $f : X \rightarrow Y$ be a morphism of complexes. If $f^k$ and $f^{k+1}$ are isomorphisms for some $k \in \left\{0, 1, \ldots, n\right\}$, then $f$ induces an isomorphism in $\Hcat{\cat{A}}$.
\end{prop}

The proof of this result can be found in \cite[Proposition~2.7]{Jasso2016Aug}.

All of this preliminary work being done, we can finally proceed to define $n$-exact categories.

\section{$n$-exact categories}

With the proper tools put in place, it is quite straightforward to expand the definition of exact categories to longer exact sequences. Similar to what was done in chapter 1, we proceed with admissible $n$-exact sequences, admissible monics, and  admissible epics:

\begin{defn}
    Let $\cat{E}$ be a class of $n$-exact sequences in $\cat{A}$. An \emph{$\cat{E}$-admissible $n$-exact sequence} is an $n$-exact sequence that is in $\cat{E}$. 
    An \emph{$\cat{E}$-admissible monic} is a morphism $d_X^{0}: X^0 \rightarrow X^1$ such that there is an $\cat{E}$-admissible $n$-exact sequence 
    \[X^0 \xra{d_X^{0}} X^1 \xra{d_X^{1}} \cdots \xra{d_X^n} X^{n+1}.\]
    
    Dually, an \emph{$\cat{E}$-admissible epic} is a morphism $d_X^n: X^n \rightarrow X^{n+1}$ such that there is an $\cat{E}$-admissible $n$-exact sequence 
    \[X^0\xra{d_X^{0}} \cdots \xra{d_X^{n-1}} X^n \xra{d_X^{n}} X^{n+1}.\]
    
    Whenever the context is clear, we will simply say admissible $n$-exact sequence (respectively admissible monic, admissible epic).
    We shall write $X^{0} \xratshit{d_{X}^{0}} X^{1}$ for admissible monics and $X^{n} \xrathshit{d_X^n} X^{n+1}$ for admissible epics.
\end{defn}

The following definition simply adapts the axioms of exact categories to the definitions given in the previous section:

\begin{defn}
    Let $\cat{E}$ be a class of $n$-exact sequences on $\cat{A}$. $\cat{E}$ is said to be an \emph{$n$-exact structure} on $\cat{A}$ if it is closed under weak isomorphisms and it satisfies the following axioms:
    \begin{sloppypar}
    \begin{itemize}[align=right]
        \item[{[E0]}] The sequence $0 \rightarrowtail 0 \rightarrow \cdots \twoheadrightarrow 0$ is an $\cat{E}$-admissible $n$-exact sequence;
        \item[{[E1]}] The composition of two $\cat{E}$-admissible monics yields an ${\cat{E}\text{-admissible monic}}$;
        \item[{[E1$^{\op}$]}] The composition of two $\cat{E}$-admissible epics yields an ${\cat{E}\text{-admissible epic}}$;
        \item[{[E2]}] If $X: X^0 \xratshit{d_X^0} \ldots \xrathshit{d_X^n} X^{n+1}$ is an $\cat{E}$-admissible $n$-exact sequence and $X^0 \xra{f^0} Y^0$ is a morphism in $\cat{A}$, there exists an $n$-pushout of $(d_X^0, \ldots, d_X^{n-1})$ along $f^0$
        % https://q.uiver.app/#q=WzAsMTEsWzAsMCwiWDoiXSxbMCwyLCJZOiJdLFsyLDAsIlheMCJdLFsyLDIsIlleMCJdLFs0LDAsIlheMSJdLFs0LDIsIlleMSJdLFs2LDAsIlxcbGRvdHMiXSxbNiwyLCJcXGxkb3RzIl0sWzgsMCwiWF5uIl0sWzgsMiwiWV5uIl0sWzEwLDAsIlhee24rMX0iXSxbMCwxLCJmIl0sWzIsNCwiZF9YXjAiLDAseyJzdHlsZSI6eyJ0YWlsIjp7Im5hbWUiOiJtb25vIn19fV0sWzMsNSwiZF9ZXnswfSIsMCx7InN0eWxlIjp7InRhaWwiOnsibmFtZSI6Im1vbm8ifSwiYm9keSI6eyJuYW1lIjoiZGFzaGVkIn19fV0sWzIsMywiZl4wIl0sWzQsNSwiZl4xIiwwLHsic3R5bGUiOnsiYm9keSI6eyJuYW1lIjoiZGFzaGVkIn19fV0sWzQsNiwiZF9YXjEiXSxbNSw3LCJkX1leezF9IiwwLHsic3R5bGUiOnsiYm9keSI6eyJuYW1lIjoiZGFzaGVkIn19fV0sWzYsOCwiZF9YXntuLTF9Il0sWzcsOSwiZF9ZXntuLTF9IiwwLHsic3R5bGUiOnsiYm9keSI6eyJuYW1lIjoiZGFzaGVkIn19fV0sWzgsOSwiZl5uIiwwLHsic3R5bGUiOnsiYm9keSI6eyJuYW1lIjoiZGFzaGVkIn19fV0sWzgsMTAsImRfWF57bn0iLDAseyJzdHlsZSI6eyJib2R5Ijp7Im5hbWUiOiJkb3R0ZWQifSwiaGVhZCI6eyJuYW1lIjoiZXBpIn19fV1d
        \[\begin{tikzcd}[ampersand replacement=\&,sep=small]
            {X:} \&\& {X^0} \&\& {X^1} \&\& \cdots \&\& {X^n} \&\& {X^{n+1}} \\
            \\
            {Y:} \&\& {Y^0} \&\& {Y^1} \&\& \cdots \&\& {Y^n}
            \arrow["f", from=1-1, to=3-1]
            \arrow["{d_X^0}", tail, from=1-3, to=1-5]
            \arrow["{f^0}", from=1-3, to=3-3]
            \arrow["{d_X^1}", from=1-5, to=1-7]
            \arrow["{f^1}", dashed, from=1-5, to=3-5]
            \arrow["{d_X^{n-1}}", from=1-7, to=1-9]
            \arrow["{d_X^{n}}", dotted, two heads, from=1-9, to=1-11]
            \arrow["{f^n}", dashed, from=1-9, to=3-9]
            \arrow["{d_Y^{0}}", dashed, tail, from=3-3, to=3-5]
            \arrow["{d_Y^{1}}", dashed, from=3-5, to=3-7]
            \arrow["{d_Y^{n-1}}", dashed, from=3-7, to=3-9]
        \end{tikzcd}\]
        such that $d_Y^0$ is an $\cat{E}$-admissible monic;
        \item[{[E2$^{\op}$]}] If $Y: Y^0 \xratshit{d_Y^0} \ldots \xrathshit{d_Y^n} Y^{n+1}$ is an $\cat{E}$-admissible $n$-exact sequence and $X^{n+1} \xra{f^{n+1}} Y^{n+1}$ is a morphism in $\cat{A}$, there is an $n$-pullback of $(d_Y^1, \ldots, d_Y^n)$ along $f^{n+1}$
        % https://q.uiver.app/#q=WzAsMTEsWzAsMCwiWDoiXSxbMCwyLCJZOiJdLFs0LDAsIlheMSJdLFs2LDAsIlxcbGRvdHMiXSxbOCwwLCJYXntufSJdLFs0LDIsIlleMSJdLFs2LDIsIlxcbGRvdHMiXSxbOCwyLCJZXntufSJdLFsxMCwyLCJZXntuKzF9Il0sWzEwLDAsIlhee24rMX0iXSxbMiwyLCJZXjAiXSxbMCwxLCJmIl0sWzIsMywiZF9YXjEiLDAseyJzdHlsZSI6eyJib2R5Ijp7Im5hbWUiOiJkYXNoZWQifX19XSxbMyw0LCJkX1hee24tMX0iLDAseyJzdHlsZSI6eyJib2R5Ijp7Im5hbWUiOiJkYXNoZWQifX19XSxbNSw2LCJkX1leMSJdLFs2LDcsImRfWV57bi0xfSJdLFs3LDgsImRfWV5uIiwwLHsic3R5bGUiOnsiaGVhZCI6eyJuYW1lIjoiZXBpIn19fV0sWzQsOSwiZF9YXntufSIsMCx7InN0eWxlIjp7ImJvZHkiOnsibmFtZSI6ImRhc2hlZCJ9LCJoZWFkIjp7Im5hbWUiOiJlcGkifX19XSxbOSw4LCJmXntuKzF9Il0sWzEwLDUsImRfWV4wIiwwLHsic3R5bGUiOnsidGFpbCI6eyJuYW1lIjoibW9ubyJ9LCJib2R5Ijp7Im5hbWUiOiJkb3R0ZWQifX19XSxbMiw1LCJmXjEiLDAseyJzdHlsZSI6eyJib2R5Ijp7Im5hbWUiOiJkYXNoZWQifX19XSxbNCw3LCJmXm4iLDEseyJzdHlsZSI6eyJib2R5Ijp7Im5hbWUiOiJkYXNoZWQifX19XV0=
        \[\begin{tikzcd}[ampersand replacement=\&,sep=small]
            {X:} \&\&\&\& {X^1} \&\& \cdots \&\& {X^{n}} \&\& {X^{n+1}} \\
            \\
            {Y:} \&\& {Y^0} \&\& {Y^1} \&\& \cdots \&\& {Y^{n}} \&\& {Y^{n+1}}
            \arrow["f", from=1-1, to=3-1]
            \arrow["{d_X^1}", dashed, from=1-5, to=1-7]
            \arrow["{f^1}", dashed, from=1-5, to=3-5]
            \arrow["{d_X^{n-1}}", dashed, from=1-7, to=1-9]
            \arrow["{d_X^{n}}", dashed, two heads, from=1-9, to=1-11]
            \arrow["{f^n}"{description}, dashed, from=1-9, to=3-9]
            \arrow["{f^{n+1}}", from=1-11, to=3-11]
            \arrow["{d_Y^0}", dotted, tail, from=3-3, to=3-5]
            \arrow["{d_Y^1}", from=3-5, to=3-7]
            \arrow["{d_Y^{n-1}}", from=3-7, to=3-9]
            \arrow["{d_Y^n}", two heads, from=3-9, to=3-11]
        \end{tikzcd}\]
        such that $d_X^n$ is an $\cat{E}$-admissible epic.
    \end{itemize}
    \end{sloppypar}
\end{defn}

\begin{rmks}
    \begin{itemize}
        \item The axiom $[E0]$ given here seems to be less restrictive than the $[E0]$ and $[E0^{op}]$ of exact categories. However, if $n \geq 2$, note that identity morphisms are also admissible monics since 

        % https://q.uiver.app/#q=WzAsMTAsWzAsMCwiWCJdLFsxLDAsIlgiXSxbMiwwLCIwIl0sWzMsMCwiXFxjZG90cyJdLFs0LDAsIjAiXSxbMCwxLCIwIl0sWzEsMSwiMCJdLFsyLDEsIjAiXSxbMywxLCJcXGNkb3RzIl0sWzQsMSwiMCJdLFsxLDJdLFsyLDNdLFszLDRdLFswLDVdLFs1LDZdLFs2LDddLFs3LDhdLFs4LDldLFsxLDZdLFsyLDddLFs0LDldLFswLDEsIjFfWCJdXQ==
        \[\begin{tikzcd}[ampersand replacement=\&]
            X \& X \& 0 \& \cdots \& 0 \\
            0 \& 0 \& 0 \& \cdots \& 0
            \arrow["{1_X}", from=1-1, to=1-2]
            \arrow[from=1-1, to=2-1]
            \arrow[from=1-2, to=1-3]
            \arrow[from=1-2, to=2-2]
            \arrow[from=1-3, to=1-4]
            \arrow[from=1-3, to=2-3]
            \arrow[from=1-4, to=1-5]
            \arrow[from=1-5, to=2-5]
            \arrow[from=2-1, to=2-2]
            \arrow[from=2-2, to=2-3]
            \arrow[from=2-3, to=2-4]
            \arrow[from=2-4, to=2-5]
        \end{tikzcd}\]
        is a weak isomorphism. Identity morphisms are admissible epics for the same reason.\\
        \item One can show that $n$-kernels and $n$-cokernels are unique up to weak isomorphism. As such, all $n$-kernels (respectively \break $n$-cokernels) of an admissible epic (respectively admissible \break monic) fit into an admissible $n$-exact sequence.
    \end{itemize}
\end{rmks}

\section{$n$-cluster tilting subcategories}
Similar to extension-closed subcategories of abelian categories, there is a higher degree notion that gives us natural examples of $n$-exact categories. But this requires a few more definitions, namely what being functorially finite means.

\begin{defn}
Let $\cat{A}$ be an additive category and $\cat{D}$ be a full subcategory of $\cat{A}$. $\cat{D}$ is called \emph{covariantly finite} if for any object $A \in \cat{A}$ there exists an object $L_A \in \cat{D}$ and a morphism $f_A : A \xra{} L_A$ such that for any morphism $g: A \xra{} D$ where $D \in \cat{D}$, there exists a morphism $h: L_A \xra{} D$ such that $g = h \circ f_A$.

Dually, $\cat{D}$ is called \emph{contravariantly finite} if for any object $A \in \cat{A}$ there exists and object $R_A \in \cat{D}$ and a morphism $f_A: R_A \xra{} A$ such that for any morphism $g: D \xra{} A$ where $D \in \cat{D}$, there exists a morphism $h: D \xra{} R_A$ such that $g = f_A \circ h$.

Both properties can be summarized in the following diagrams:

    % https://q.uiver.app/#q=WzAsMTAsWzAsMl0sWzEsMl0sWzAsMCwiQSJdLFsxLDAsIkxfQSJdLFsxLDEsIkQiXSxbMywyXSxbNCwyXSxbMywwLCJSX0EiXSxbNCwwLCJBIl0sWzMsMSwiRCJdLFswLDEsIlxcdGV4dHtjb3ZhcmlhbnRseSBmaW5pdGV9IiwxXSxbMiwzLCJmX0EiXSxbMyw0LCJcXGV4aXN0cyBnIiwwLHsic3R5bGUiOnsiYm9keSI6eyJuYW1lIjoiZGFzaGVkIn19fV0sWzIsNCwiaCIsMl0sWzUsNiwiXFx0ZXh0e2NvbnRyYXZhcmlhbnRseSBmaW5pdGV9IiwxLHsic3R5bGUiOnsiYm9keSI6eyJuYW1lIjoibm9uZSJ9LCJoZWFkIjp7Im5hbWUiOiJub25lIn19fV0sWzcsOCwiZl9BIl0sWzksOCwiaCIsMl0sWzksNywiXFxleGlzdHMgZyIsMCx7InN0eWxlIjp7ImJvZHkiOnsibmFtZSI6ImRhc2hlZCJ9fX1dXQ==
    \[\begin{tikzcd}[ampersand replacement=\&]
        A \& {L_A} \&\& {R_A} \& A \\
        \& D \&\& D \\
        {} \& {} \&\& {} \& {}
        \arrow["{f_A}", from=1-1, to=1-2]
        \arrow["h"', from=1-1, to=2-2]
        \arrow["{\exists g}", dashed, from=1-2, to=2-2]
        \arrow["{f_A}", from=1-4, to=1-5]
        \arrow["{\exists g}", dashed, from=2-4, to=1-4]
        \arrow["h"', from=2-4, to=1-5]
        \arrow["{\text{covariantly finite}}"{description}, from=3-1, to=3-2]
        \arrow["{\text{contravariantly finite}}"{description}, draw=none, from=3-4, to=3-5]
    \end{tikzcd}\]

    $\cat{D}$ is then said to be \emph{functorially finite} if it is both covariantly and contravariantly finite.
\end{defn}

\begin{rmk}
In the context of module categories, this is very often true for any subcategory. In fact, if $A$ is an Artin algebra, one can show that every subcategory of $\text{mod}\ A$ is functorially finite if and only if $A$ is of finite representation type (see \cite[Proposition~1.2]{AUSLANDER1991111}).
\end{rmk}



\begin{defn}
Let $\cat{A}$ be an abelian category and $\cat{D}$ be a full subcategory of $\cat{A}$. $\cat{D}$ is an \emph{$n$-cluster tilting category} if
\begin{itemize}
    \item it is functorially finite;
    \item $\cat{D} = \left\{C \in \cat{A} \ \middle|\  \text{Ext}_{\cat{A}}^i(C, D) = 0 \ \forall D \in \cat{D}\right\}$;
    \item $\cat{D} = \left\{C \in \cat{A} \ \middle|\  \text{Ext}_{\cat{A}}^i(D, C) = 0 \ \forall D \in \cat{D}\right\}$.
\end{itemize}
\end{defn}